%============================================================================
% prelude.tex
%   - titlepage
%   - table of contents, list of tables and list of figures
%   - abstract
%============================================================================
\clearpage%
\pagenumbering{roman}  % This makes the page numbers Roman (i, ii, etc)
%--------------------------------------------------------------------%
% TITLE PAGE
%   - define \title{} \author{} \date{}

\title{ Title of the project  }
\author{Student Name (13EL1001)}

\date{Oct 2015}

\department{Department of Electronics and Telecommunication Engineering}

%    - Set the guide's name
\setguide{Dr. Mukesh D. Patil}
%    - Set the coguide's name (if you have one)
%\setcoguide{Prof. Name of the Guide}

%   - once the above are defined, use \maketitle to generate the titlepage
\maketitle

%--------------------------------------------------------------------
\cleardoublepage \nonumber
\begin{center}
\includegraphics{raitlogo.eps}\\
Ramrao Adik Education Society's\\
\large {\textbf{Ramrao Adik Institute of Technology}}\\
\small(Affiliated to the University of Mumbai)\\
Dr. D. Y. Patil Vidyanagar,Sector 7, Nerul, Navi Mumbai 400
706.\\\vspace{0.3in}
\textbf{\begin{Huge}Certificate\end{Huge}}
\end{center}
\vspace{0.1in}
\begin{center}
\large This is to certify that, the special topic seminar entitled
\vspace{0.15in}\\\textbf {\lq\lq Visible light Communication and its Applications\rq\rq} \\
\vspace{0.1in}is a bonafide work done by\vspace{0.15in}
\\\textbf { Student Name} \\
(Roll No. 13EL1001)\\\vspace{0.1in} and is submitted in the partial
fulfillment of
the requirement for the degree of \vspace{0.15in} \\\textbf{Bachelors of Engineering}\\
in \\\textbf{Electronics AND Telecommunication Engineering} \\to the \\\textbf{University
of Mumbai}.


\includegraphics{university.eps}

 \vspace{0.2in}
%
%\begin{tabular}{cccc}
%      \rule{5cm}{1sp}
%      \\
%       Guide \\\\
%      \rule{5cm}{1sp}  & \rule{5cm}{1sp}   & \rule{5cm}{1sp} \\
%      Project Coordinator      & Head of Department         & Principal \\
%
%
%    \end{tabular}
%\vspace{1in}
%%%%%%%%%%%%%%%%%%%%%%%%%%%%%%%%%%%%%%%%%%%%%%%%%%%%%%%%%%%%%%%%%%%%%%%%%%
\begin{tabular}{ccc}

      \rule{4.0cm}{1sp} & \hspace{0.7in}                & \hspace{0.3in} \rule{4.0cm}{1sp} \\ \vspace{0.65in}
       Examiner & \hspace{1in}             & \hspace{0.3in} Supervisor \\
      \rule{4.0cm}{1sp} & \hspace{0.5in} \rule{4.0cm}{1sp}                  &\hspace{0.45in} \rule{4.0cm}{1sp} \\
      Project Coordinator & \hspace{0.5in} Head of Department                &\hspace{0.65in} Principal \\\\
        %&\rule{5cm}{1sp}& \\
%       &Principal&
    \end{tabular}
\end{center}




%%%%%%%%%%%%%%%%%%%%%%%%%%%%%%%%%%%%%%%%%%%%%%%%%%%%%%%%%%%%%%%%%%%%%%%%%%%%%%%%%%%%%%%%%%%%%%%%%%
%\begin{tabular}{ccc}
%
%      \rule{5cm}{1sp}                && \rule{5cm}{1sp} \\ \vspace{0.65in}
%       Guide             && Project Coordinator \\
%      \rule{5cm}{1sp}                && \rule{5cm}{1sp} \\
%      Head of Department                && Principal \\\\
%        %&\rule{5cm}{1sp}& \\
%%       &Principal&
%    \end{tabular}
%\end{center}
%-----------------------------------------------------------------
\cleardoublepage
%--------------------------------------------------------------------%

%\centerline{\textbf{\begin{Huge}Certificate of Approval by
%Examiners\end{Huge}}}
%
%\vspace{1.5in}
%
% This is to certify that the Stage-1 report entitled
%"\textbf{Title of the project}" is a bonafide work done by
%\textbf{Name of the student} under the guidance of \textbf{Name of
%the guide}. This dissertation has been approved for the award of
%\textbf{Bachelor's Degree in Electronics  Engineering}, University of
%Mumbai.
%
%\vspace{1.5in} \textbf{Examiners:}\vspace{1in}\\
%
%\begin{tabular}{ccc}
%
%      \rule{5cm}{1sp} &\rule{30mm}{0pt} & \rule{5cm}{1sp}
%      \\\vspace{0.1in}
%       Name:                && Name:  \\
%
%    \end{tabular}
%



\begin{acknowledgments}
I would like to express my deepest gratitude to my guide, Dr. Mukesh D. Patil, for
his excellent guidance caring, patience, and providing me with an excellent atmosphere which helped me in successful completion of this special topic.
%for doing this .
%With great pleasure, I avail this opportunity to express my deep
%sense of gratitude to my guide, Dr. Mukesh D. Patil, for his spirited guidance
%and inspiration.I have deep sense of admiration for their innate
%goodness and inexhaustible enthusiasm.it helped me to work in right
%direction to attain desired objective.

I am also thankful to our MiniProject-I Coordinator, Mrs. Ashwini Joshi who devoted their
valuable time and helped me in all possible ways towards successful
completion of this special topic. I am also grateful to all those who have contributed
directly or indirectly to this work.

I am thankful to our Principal Dr. Ramesh Vasappanavara  for their support and encouragement.
%I extend thanks to my friends who have done lots of
%nice things for me. I cannot end without thanking my lovely family
%for their encouragement.

I would also like to express my deepest thanks to all my friends and my colleagues, who provided a perfect environment for my studies.

Last but not least, I would like to express my heartily gratitudes to my family for
their unconditional support and love.
\end{acknowledgments}

% ABSTRACT
\begin{abstract}

Visible light communications (VLC) is a technology for wireless communication using light that can be perceived by the naked eye. VLC uses frequencies other than radio, and they are unrestricted and licence free. The urgent need of VLC is to overcome the problems faced in RF communication. Unlike existing methods of wireless communication, the visible light portion of the electromagnetic frequency spectrum is used in VLC  to transmit information. VLC refers to the communication technology which utilizes the visible light source as a signal transmitter, the air as a transmission medium, and the appropriate photodiode as a signal receiving component. This report provides  an overview of applications and design challenges for VLC, compare it with other existing communication technologies and presents the modulation techniques used.


\end{abstract}

%--------------------------------------------------------------------%
% CONTENTS, TABLES, FIGURES
 \listoftables
 \listoffigures
 \tableofcontents

%--------------------------------------------------------------------%
%  Single counter for theorems and theorem-like environments:
\newtheorem{theorem}{Theorem}[chapter]
\newtheorem{assertion}[theorem]{Assertion}
\newtheorem{claim}[theorem]{Claim}
\newtheorem{conjecture}[theorem]{Conjecture}
\newtheorem{corollary}[theorem]{Corollary}
\newtheorem{definition}[theorem]{Definition}
\newtheorem{example}[theorem]{Example}
\newtheorem{figger}[theorem]{Figure}
\newtheorem{lemma}[theorem]{Lemma}
\newtheorem{prop}[theorem]{Proposition}
\newtheorem{remark}[theorem]{Remark}

%--------------------------------------------------------------------%
% Make the page numbers Arabic (1, 2, etc)
\cleardoublepage%
\pagenumbering{arabic} 